% Finding A15. Prepared independently of final/.  One-fold Frobenius twist
% notation C^{(1)} agrees with the root paper's revised convention.
\subsection{Rank one permits a primitive factor of arbitrary odd conductor}
\label{sec:rank-one-composite-bridge}

\paragraph{Finding A15 (the composite primitive-factor criterion).}
Let $C/k$ be a smooth projective geometrically connected curve of genus
two over a finite field of characteristic two, with geometric $2$-rank
$f_C=1$. Let $n>1$ be odd. The following conditions suffice for the
point-section separation conclusion of A1.
First, there is a geometrically connected cyclic etale cover
$D\to C$ of degree $n$ whose primitive conductor-$n$ Jacobian factor
has $2$-rank zero. More generally, for a specified $P\in C(k)$ it
suffices to have a finite extension $L/k$, such a cover over $L$ with
constant deck group $G=\langle\sigma\rangle$, and a point
$P'\in D(L)$ over $P$, for which
\[
 A=\prod_{\ell\mid n}(1-\sigma^{n/\ell}),\qquad
 P_n=\operatorname{im}(A:J_D\longrightarrow J_D),\qquad
 P_n(L)[2^\infty]=0.                                      \tag{CB.1}
\]
The product is over distinct prime divisors of $n$.
In particular, the first condition holds if the Raynaud theta divisor
on $\operatorname{Pic}^0(C^{(1)})$ contains a point of exact odd order
$n$ and
\[
                    \langle 2,-1\rangle=(\mathbf Z/n)^\times.
                                                               \tag{CB.2}
\]
Here $C^{(1)}$ is the one-fold relative Frobenius twist. Thus for
$2$-rank one the order $n$ need not be a prime power. No assertion of
universal theta incidence is included.

\paragraph{Exact source scope.}
The geometric criterion underlying (CB.1), with the extra hypothesis
that the base is not geometrically a separable double cover of a
genus-one curve, is the recovered \texttt{ANABELIAN\_SUCCESSOR}
\emph{final row A.14}, working identifier AC.11, in
\path{agent_criteria/criteria_section.tex}. The refinement here
removes that extra hypothesis for the rank-one genus-two stratum and
combines the criterion with theta incidence. Final row A.35 is a
different, finite Fourier counting criterion; it is not the source of
the composite projector argument. We give the full argument below.

\begin{proof}
\emph{1. Rank one excludes the exceptional double quotient.}
Work first over an algebraic closure. Suppose that a genus-two curve
$C$ admits a separable degree-two map to a genus-one curve $E$.
The function-field extension is quadratic and Galois, with group
$\mathbf Z/2$. Riemann--Hurwitz says that its different has degree
$2g_C-2=2$. Every ramified point in this characteristic-two quadratic
extension has different exponent at least two, so there is exactly
one ramified point. The Deuring--Shafarevich formula consequently gives
\[
 f_C-1=2(f_E-1)+1,\qquad f_C=2f_E\in\{0,2\}.              \tag{CB.3}
\]
This contradicts $f_C=1$. We use the geometric $p$-rank form of that
formula, stated for $p$-group covers, for example, in
Anbar--Beelen--Nguyen, Theorem~1, equation~(1)
\cite{AnbarBeelenNguyenP}; no assumption of ordinarity is needed.

\emph{2. The primitive projector produces a small separable function.}
Suppose that distinct points $P,Q$ collided in the covering formulation
of A1. The pointed cover in (CB.1), after restricting the collision
to $L$, supplies $Q'\in D(L)$ above $Q$ with
$\delta=[Q'-P']\in J_D(L)[2^\infty]$. Hence $A\delta=0$.
Put $r=\omega(n)$ and $d=2^r$. Each element
$\sigma^{n/\ell}$ has order $\ell$, and these elements belong to
distinct prime-primary factors of the cyclic group. Their $2^r$
subset products are therefore all distinct. Expanding $A(Q'-P')$
gives a nonzero principal divisor with exactly $d$ simple zeros and
$d$ simple poles. They lie in the two disjoint free deck orbits of
$P'$ and $Q'$. Let $f$ be the resulting rational function on
$D_{\overline L}$. It has degree $d$ and is separable: its simple
zeros (or poles) rule out a square in the function field.

The relevant degree bound is
\[
                         n>(2^{\omega(n)}-1)^2=(d-1)^2.
                                                               \tag{CB.4}
\]
It holds for every odd $n>1$, with no additional size assumption.
For $r=1$ use $n\ge3>1$; for $r=2$ use $n\ge15>9$.
For $r\ge3$, the distinct odd prime factors give
$n\ge3\cdot5^{r-1}>4^r>(2^r-1)^2$; the middle inequality holds
at $r=3$ and its ratio increases by $5/4$ thereafter.

\emph{3. The field generated by the conjugate functions.}
Write $K=\overline L(D)$,
$M=\overline L(\gamma f:\gamma\in G)$, and $e=[K:M]$.
Poonen's Castelnuovo--Severi consequence gives
\[
                  e\mid d,\qquad g(M)\le(d/e-1)^2.
                                                               \tag{CB.5}
\]
Here is the entire induction needed from that consequence. Start with
one degree-$d$ function, whose rational field has index $d$ and genus
zero. If an intermediate field has index $d/s$ and genus at most
$(s-1)^2$, adjoining another degree-$d$ function with relative degree
$a$ gives index $d/(as)$. Castelnuovo--Severi bounds its genus by
\[
 a(s-1)^2+(a-1)(as-1)
  =(as-1)^2-as(a-1)(s-1)\le(as-1)^2.
\]
Induction over the finitely many conjugates proves (CB.5).
This is exactly the use of Proposition~2.1, Lemma~2.2, and
Corollary~2.3 of \cite{Poonen2007Gonality}, over the perfect field
$\overline L$.

Since $D\to C$ is etale, $g(D)=n+1$. Equations (CB.4)--(CB.5)
exclude $e=1$. Thus $e>1$ is a power of two. Let $Z$ be the smooth
projective curve with function field $M$, and let $v:D\to Z$ have
degree $e$. The group $G$ preserves $M$. Its induced action on $Z$
is faithful, since a kernel $H$ would satisfy
$|H|\mid[K:M]=e$ as well as $|H|\mid n$.
The action on $Z$ is also free. Indeed, a stabilizer of a geometric
point of $Z$ would contain an element of odd prime order $\ell$.
It would partition the fibre divisor of $v$ into orbits of length
$\ell$, since the deck action on $D$ is free; fibre multiplicities
are constant on these orbits. This forces $\ell\mid e$, impossible.

Riemann--Hurwitz for the free action now yields
\[
                    g(Z)-1=n\bigl(g(Z/G)-1\bigr).
\]
The bound $0\le g(Z)\le(d/e-1)^2<n$ forces
$g(Z)=g(Z/G)=1$: quotient genus zero would give $g(Z)=1-n<0$,
and quotient genus at least two would give $g(Z)\ge n+1$.

\emph{4. A reduced fibre forces the forbidden degree-two map.}
Write $f=h\circ v$, with $h$ a nonconstant function on $Z$.
Choose a zero $z$ of $h$. The fibre divisor $v^*z$ is reduced,
because the zero divisor of $f$ is reduced. It has $e$ distinct
geometric points, all contained in the two original deck orbits.
Each such orbit maps injectively to $Z$, by equivariance and the
free action on $Z$. Hence the fibre contains at most one point from
each orbit, giving $e\le2$. Thus $e=2$.

The extension $K/\overline L(f)$ is separable, so $K/M$ is separable.
Its quadratic involution $\tau$ is normalized by $G$, because $G$
preserves its fixed field $M$; it is centralized by $G$ since
$\operatorname{Aut}(\mathbf Z/2)=1$. Therefore taking the quotient
by $G$ yields a separable degree-two map
\[
                  C=D/G\longrightarrow Z/G
\]
to a genus-one curve. Its degree is two explicitly because
$[K:M^G]=2n$ and $[K:K^G]=n$.
This contradicts Step~1 and proves the
criterion (CB.1). A primitive factor of geometric $2$-rank zero
has no nonzero geometric $2$-power torsion points. Its rational
$2$-primary group is therefore zero after every finite constant
extension, so this geometric condition implies the required
pointed version of (CB.1) for every base point.

\emph{5. Theta incidence annihilates every primitive orbit.}
Let $L$ have exact order $n$ on the Raynaud theta divisor of
$C^{(1)}$, and let $B$ denote its locally exact differential bundle.
In characteristic two $B$ is a line bundle with $B^2\simeq\omega$.
Serre duality therefore gives
$h^0(B\otimes L)=h^0(B\otimes L^{-1})$, so both signs lie on theta.
For relative Frobenius $F:C\to C^{(1)}$, tensor its exact sequence
by $L$. The induced map
\[
 H^1(C^{(1)},L)\longrightarrow H^1(C,F^*L)
\]
has a nonzero kernel, since $H^0(B\otimes L)\ne0$.
Both spaces are lines: the two line bundles are nontrivial of
degree zero on genus-two curves. The map is consequently zero.
The same holds for $L^{-1}$.

Choose $b$ with $2b\equiv1\pmod n$ and put $M_0=(F^*L)^b$ on $C$.
Then $M_0$ has exact order $n$ and
$F^*(M_0^{(1)})=M_0^2=F^*L$. Since Frobenius pullback is
injective on odd-order torsion, $M_0^{(1)}=L$. On the connected
cyclic cover associated to $M_0$, label the eigensheaves by $M_0^i$.
For $i\ne0\bmod n$ the cohomology spaces
$V_i=H^1(C,M_0^i)$ are one-dimensional, and Frobenius maps
$V_i$ to $V_{2i}$. The arrows $1\to2$ and $-1\to-2$ just shown
are zero. Condition (CB.2) makes every doubling orbit in
$(\mathbf Z/n)^\times$ contain $1$ or $-1$. Thus every primitive
orbit has nilpotent Frobenius, and its stable part is zero.

Finally $A$ acts on character $i$ by
$\prod_{\ell\mid n}(1-\zeta_\ell^i)$, which is nonzero precisely
when $\gcd(i,n)=1$. Since $n$ is odd, the character projectors are
defined after an extension of $\mathbf Z_2$ and split the
$2$-primary representation. The image $P_n$ is exactly the
primitive factor up to isogeny; its stable Frobenius part is the
primitive sum just considered. It therefore has $2$-rank zero.
All the geometric cover data descend after a finite constant
extension, proving the stated theta consequence.
\end{proof}

A18 below removes this geometric exclusion for the geometric
primitive-rank-zero criterion in characteristic two, in every genus
at least two. The rational-only criterion retains the distinction.

Steps~2--5 apply to any genus-two base known not to admit a geometric
separable double cover of a genus-one curve; the rank-one hypothesis
is what makes this geometric exclusion automatic in Step~1.

\paragraph{How many prime factors can the single-point condition allow?}
For odd $n$, the Chinese remainder theorem and the cyclicity of units
modulo an odd prime power give
\[
 (\mathbf Z/n)^\times/\bigl((\mathbf Z/n)^\times\bigr)^2
                  \simeq (\mathbf Z/2)^{\omega(n)}.
\]
If (CB.2) holds, this quotient is generated by the images of $2$ and
$-1$, so necessarily $\omega(n)\le2$. This is an arithmetic
restriction on the one-point sufficient condition, not on the
geometric criterion (CB.1). Exact finite examples are
\[
\begin{array}{c|rrrrrr}
 n&15&21&35&45&33&105\\ \hline
 \varphi(n)&8&12&24&24&20&48\\
 \operatorname{ord}_n(2)&4&6&12&12&10&12\\
 |\langle2,-1\rangle|&8&12&24&24&10&24
\end{array}
\]
Thus the first four integers satisfy (CB.2); having at most two
distinct prime factors is not by itself sufficient, as $n=33$ shows.
These computations certify only the unit groups. Theta incidence at
these orders is a separate condition, illustrated by A16 for $n=15$.

\paragraph{Several theta points: the exact coverage condition.}
For any fixed cyclic torsion class $M_0$ of odd order $n$, define
\[
 Z(M_0)=\{i\in(\mathbf Z/n)^\times:
                    M_0^{(1)i}\in\Theta_C\}.
\]
The primitive factor of its cyclic cover has $2$-rank zero if and
only if $Z(M_0)$ meets every orbit of multiplication by $2$ on
$(\mathbf Z/n)^\times$. Indeed, the arrow $V_i\to V_{2i}$ is
zero exactly at the displayed theta incidence. A cyclic sequence
of maps of lines has zero stable part exactly when at least one
arrow is zero; if every arrow is nonzero, their composition is a
bijective semilinear map of a line. This proves the equivalence.
Consequently, on the rank-one base, sufficiently many theta points
\emph{in this one cyclic subgroup} suffice, even when $\omega(n)>2$.
Condition (CB.2) and symmetry give this coverage from a single point.

\paragraph{What the primary-projection theorem still leaves open.}
Applying Poonen's theorem with $S_0\supset\{2,\ell\}$ can prescribe
an $\ell^a$ component, but gives order $n=\ell^a m$, where all primes
of $m$ lie outside a subsequently chosen density-one set $S$.
In the proof, Lemma~6.4 first finds a point $h$ in a bounded-index
subgroup supported on a density-zero prime set and then takes
$U$ to be \emph{all} prime divisors of its order. It bounds the
density quantity $\eta(U)$, not $|U|$ or $\omega(\operatorname{ord}h)$.
Enlarging finite $S_0$ therefore does not give the extra restriction
$\omega(m)\le1$ needed even to qualify for (CB.2), or the unit-group
equality itself \cite[Theorem~1.8 and Lemmas~6.4--6.6]{Poonen2005}.
For $\ell=3$, A10 forces $m>1$ on the rank-one stratum; A15 now
permits such a cofactor, but does not manufacture a suitable one.
Theta points of pairwise coprime orders likewise do not establish
the displayed orbit-coverage condition: on their fibre-product
cover they control proper-conductor characters, while the primitive
characters are tensor products whose theta incidence is not supplied.
This identifies the additional incidence statement that would be
needed to complete this proposed route.
